\documentclass[11pt,a4paper]{article}
\usepackage{fontspec,amsmath,amssymb,amsthm,microtype}
\setmainfont{lmroman10-regular.otf}[BoldFont=lmroman10-bold.otf,ItalicFont=lmroman10-italic.otf,BoldItalicFont=lmroman10-bolditalic.otf]
\usepackage{xcolor}
\usepackage[most]{tcolorbox}
\usepackage{mathrsfs}
\providecommand{\lt}{<}
\providecommand{\gt}{>}
\setmonofont{DejaVuSansMono.ttf}[BoldFont=DejaVuSansMono-Bold.ttf,ItalicFont=DejaVuSansMono-Oblique.ttf,BoldItalicFont=DejaVuSansMono-BoldOblique.ttf,Scale=0.86]
\newfontfamily\leanfallback{STIXTwoMath-Regular.otf}[Scale=0.86]
\definecolor{leanrule}{gray}{0.62}
\definecolor{leanback}{gray}{0.962}
\definecolor{leantitle}{gray}{0.22}
\definecolor{leanmuted}{gray}{0.38}
\newtcolorbox{leanbox}[1]{enhanced,breakable,colback=leanback,colframe=leanrule,
  boxrule=0.4pt,arc=0.6mm,left=2.4mm,right=2.4mm,top=1.8mm,bottom=1.8mm,
  before skip=0.8em,after skip=1em,fonttitle=\small\bfseries,coltitle=leantitle,
  colbacktitle=leanback,titlerule=0.3pt,title={#1}}
\newlength{\leancharwidth}
\newenvironment{leancode}{\par\smallskip\begingroup\ttfamily\footnotesize\raggedright
  \settowidth{\leancharwidth}{M}\setlength{\parskip}{0pt}\setlength{\parindent}{0pt}
  \frenchspacing}{\par\endgroup\smallskip}
\newcommand{\leanline}[2]{\hangindent=\dimexpr#1\leancharwidth+4\leancharwidth\relax
  \hangafter=1\noindent\hspace*{#1\leancharwidth}#2\par}
\newcommand{\leansource}[1]{\par{\footnotesize\color{leanmuted}#1}\par}
\newcommand{\leanlabel}[1]{\par\smallskip\noindent{\small\bfseries #1}\par}
\usepackage[margin=25mm]{geometry}
\usepackage[unicode,hidelinks]{hyperref}
\hypersetup{pdfauthor={LangevinLean Project}}
\setlength{\emergencystretch}{3em}
\setlength{\parskip}{0.4em}
\setlength{\parindent}{0pt}
\allowdisplaybreaks[2]
\raggedbottom
\title{\texorpdfstring{The strongly convex IPLA parameter-error bound}{The strongly convex IPLA parameter-error bound}}
\author{LangevinLean Project}
\date{}
\begin{document}\maketitle
Draft mathematical exposition; these notes may contain mistakes.


The parameter error is bounded by equilibrium concentration, the remaining initial transport error, and the discretization error. For the observed-data potential \(U_y\), assume differentiability, \(m\)-strong convexity with \(m>0\), an \(L\)-Lipschitz gradient, \(N>0\), and \(0<\gamma\le2/(m+L)\). Assume square-integrable initialization and mutually independent standard Gaussian innovation coordinates, independent of initialization. Let \(\theta_*\) maximize the latent-integrated likelihood. The bound below holds for arbitrary nonnegative parameter and latent dimensions.



Earlier results used here: \cite[Section~3.50]{LangevinLean2026ULASC}; \cite[Section~3.52]{LangevinLean2026ULASC}; \cite[Section~3.53]{LangevinLean2026ULASC}.

\section{\texorpdfstring{The transport and equilibrium contributions to parameter error}{The transport and equilibrium contributions to parameter error}}


Suppose the state dimension is positive. The summed potential has curvature \(Nm\) and gradient constant \(NL\) in dimension \(D=d_\theta+Nd_x\). Its Gibbs law \(\Pi_N\) is a probability with finite second moment, so square-integrable initialization gives finite initial quadratic transport. Apply Gaussian ULA with \(h=\gamma/N\). The scaled step restriction and the exact law identity give

\[W_2(\rho_k,\Pi_N)\le(1-m\gamma)^kW_2(\mu_0,\Pi_N)+\frac{33}{20}\frac Lm\sqrt{d_x+d_\theta/N}\sqrt\gamma.\]

The estimate is for a finite transport distance between the algorithm law and the Gibbs law. The equilibrium bound is \(\operatorname{RMS}_{\Pi_N}(p_\theta-\theta_*)\le\sqrt{2d_\theta/(Nm)}\). For any coupling of \(\rho_k\) and \(\Pi_N\), the projection norm bound \(\|p_\theta\|\le1\) and Minkowski’s inequality bound the parameter error by the coupling cost plus that equilibrium radius. Taking the infimum over couplings yields

\[\left(\int\|p_\theta z-\theta_*\|^2d\rho_k\right)^{1/2}\le\sqrt{\frac{2d_\theta}{Nm}}+(1-m\gamma)^kW_2(\mu_0,\Pi_N)+\frac{33}{20}\frac Lm\sqrt{d_x+\frac{d_\theta}{N}}\sqrt\gamma.\]

The moment assumptions needed in this argument follow from the joint potential hypotheses and \(L^2\) initialization. The substitutions in the Gaussian theorem are step \(\gamma/N\), curvature \(Nm\), gradient constant \(NL\), and dimension \(d_\theta+Nd_x\). For a coupling \((X,Y)\) with marginals \(\rho_k,\Pi_N\), the projection inequality used above is exactly \(\|p_\theta X-\theta_*\|_2\le\|X-Y\|_2+\|p_\theta Y-\theta_*\|_2\). Thus transport, concentration, and projection give the three displayed terms. Stationarity of the reference diffusion and the target concentration are derived results, rather than extra assumptions of this theorem.

\begin{leanbox}{Lean: \texttt{nonempty\_\allowbreak{}dimension\_\allowbreak{}bound}}
\leansource{\texttt{nonempty\_\allowbreak{}dimension\_\allowbreak{}bound} \quad IPLA/\allowbreak{}Theorem.lean, lines 41–99}
\begin{leancode}
\leanline{0}{private\ theorem\ nonempty\_dimension\_bound}
\leanline{4}{\{Ω\ :\ Type*\}\ [MeasurableSpace\ Ω]\ (P\ :\ Measure\ Ω)\ [IsProbabilityMeasure\ P]}
\leanline{4}{(N\ dθ\ dx\ :\ ℕ)\ [Nonempty\ (Fin\ dθ\ ⊕\ (Fin\ N\ ×\ Fin\ dx))]\ (hN\ :\ 0\ <\ N)}
\leanline{4}{\{Data\ :\ Type*\}\ (U\ :\ Data\ →\ Joint\ dθ\ dx\ →\ ℝ)\ (y\ :\ Data)}
\leanline{4}{(m\ γ\ :\ ℝ)\ (L\ :\ ℝ≥0)}
\leanline{4}{(hdiff\ :\ Differentiable\ ℝ\ (U\ y))\ (hsc\ :\ StrongConvexOn\ Set.univ\ m\ (U\ y))}
\leanline{4}{(hLip\ :\ LipschitzWith\ L\ (gradient\ (U\ y)))}
\leanline{4}{(hm\ :\ 0\ <\ m)}
\leanline{4}{(hγ\ :\ 0\ <\ γ)\ (hsmall\ :\ γ\ ≤\ 2\ /\ (m\ +\ (L\ :\ ℝ)))}
\leanline{4}{(θStar\ :\ EuclideanSpace\ ℝ\ (Fin\ dθ))\ (hMaximizer\ :\ IsMarginalMaximizer\ (U\ y)\ θStar)}
\leanline{4}{(initial\ :\ Lp\ (State\ N\ dθ\ dx)\ 2\ P)\ (noise\ :\ ℕ\ →\ Lp\ (State\ N\ dθ\ dx)\ 2\ P)}
\leanline{4}{(hGaussian\ :\ GaussianInnovations\ P\ initial\ noise)}
\leanline{4}{(k\ :\ ℕ)\ :}
\leanline{4}{Real.sqrt\ (∫\ z,\ ‖parameterCLM\ z\ -\ θStar‖\ \textasciicircum{}\ 2}
\leanline{6}{∂(laws\ P\ (U\ y)\ γ\ L\ hdiff\ hLip\ initial\ noise\ k))\ ≤}
\leanline{6}{Real.sqrt\ (2\ *\ (dθ\ :\ ℝ)\ /\ ((N\ :\ ℝ)\ *\ m))\ +}
\leanline{8}{(1\ -\ m\ *\ γ)\ \textasciicircum{}\ k\ *\ ULA.wasserstein2Real\ (P.map\ initial)\ (gibbsTarget\ (U\ y))\ +}
\leanline{8}{(33\ /\ 20\ :\ ℝ)\ *\ ((L\ :\ ℝ)\ /\ m)\ *}
\leanline{10}{Real.sqrt\ ((dx\ :\ ℝ)\ +\ (dθ\ :\ ℝ)\ /\ (N\ :\ ℝ))\ *\ Real.sqrt\ γ}
\end{leancode}
\leanlabel{How the Lean proof follows the argument}
\texttt{nonempty\_\allowbreak{}dimension\_\allowbreak{}bound} first derives target normalization and \texttt{hfinite}, the finite initial transport. \texttt{hs} is the rescaled step restriction. \texttt{htransport} applies \texttt{ULA.\allowbreak{}Hilbert.\allowbreak{}dalalyan\_\allowbreak{}of\_\allowbreak{}gaussian\_\allowbreak{}innovations} to the summed potential, passing each field of \texttt{h\allowbreak{}Gaussian}. \texttt{houtput} separately proves finite output transport, and \texttt{hlaw} identifies its numerical law. \texttt{hfactor} and \texttt{scaled\_\allowbreak{}bias\_\allowbreak{}eq} cancel particle factors. Finally \texttt{hconcentration}, \texttt{hprojected} and \texttt{hmoment} give the finite parameter root moment. Rewriting the projected norm as its integral and adding the transport estimate closes the proof.



The projection inequality gives \(R_{\theta_*}(\rho_k)\le W_2(\rho_k,\Pi_N)+R_{\theta_*}(\Pi_N)\); equilibrium concentration supplies the second term. The facts \texttt{houtput} and \texttt{hmoment} ensure finite values before projected RMS is rewritten as a square-root integral. Finally, the exact particle/\allowbreak{}time scaling simplifies the transport terms. Thus normalization, stationary reference and parameter concentration are derived inputs, while joint strong convexity, smoothness and the Gaussian innovation laws remain the stated assumptions.



Formal correspondence: nonempty\_dimension\_bound : For fixed observations \(y\), write \(U_y(\operatorname{pair}(\theta,x))=-\log p_\theta(x,y)\). The general-dimension counterpart of equation (34), with independent parameter dimension \(d_\theta\) and latent dimension \(d_x\), for the actual Gaussian IPLA recursion and its defined Gibbs target. Every hypothesis is displayed here or is one of the Gaussian-law and independence fields of GaussianInnovations , whose definition is included in the standalone file. Normalization, all finite transport moments, a stationary Langevin reference, and the equilibrium parameter concentration are derived in the proof.
\end{leanbox}


Earlier results used here: \cite[Section~3.17.4]{LangevinLean2026ULASC}; \cite[Section~3.26.3]{LangevinLean2026ULASC}; \cite[Section~3.45.6]{LangevinLean2026ULASC}; \cite[Section~3.48.6]{LangevinLean2026ULASC}; \cite[Section~3.48.7]{LangevinLean2026ULASC}; \cite[Section~3.48.9]{LangevinLean2026ULASC}; \cite[Section~3.50.2]{LangevinLean2026ULASC}; \cite[Section~3.50.3]{LangevinLean2026ULASC}.

\section{\texorpdfstring{The zero-dimensional cases}{The zero-dimensional cases}}


Under \(N>0\), a positive parameter dimension supplies a parameter coordinate, and a positive latent dimension supplies a coordinate in a nonempty particle block. Hence the full state has no coordinates exactly when both dimensions are zero. Every other case is covered by the positive-state-dimension argument. In the remaining singleton case the parameter error, transport distance, and dimension-dependent error terms are all zero, so the same bound holds directly.

\begin{leanbox}{Lean: \texttt{theorem1}}
\leansource{\texttt{theorem1} \quad IPLA/\allowbreak{}Theorem.lean, lines 105–137}
\begin{leancode}
\leanline{0}{theorem\ theorem1}
\leanline{4}{\{Ω\ :\ Type*\}\ [MeasurableSpace\ Ω]\ (P\ :\ Measure\ Ω)\ [IsProbabilityMeasure\ P]}
\leanline{4}{(N\ dθ\ dx\ :\ ℕ)\ (hN\ :\ 0\ <\ N)}
\leanline{4}{\{Data\ :\ Type*\}\ (U\ :\ Data\ →\ Joint\ dθ\ dx\ →\ ℝ)\ (y\ :\ Data)}
\leanline{4}{(m\ γ\ :\ ℝ)\ (L\ :\ ℝ≥0)}
\leanline{4}{(hdiff\ :\ Differentiable\ ℝ\ (U\ y))\ (hsc\ :\ StrongConvexOn\ Set.univ\ m\ (U\ y))}
\leanline{4}{(hLip\ :\ LipschitzWith\ L\ (gradient\ (U\ y)))}
\leanline{4}{(hm\ :\ 0\ <\ m)}
\leanline{4}{(hγ\ :\ 0\ <\ γ)\ (hsmall\ :\ γ\ ≤\ 2\ /\ (m\ +\ (L\ :\ ℝ)))}
\leanline{4}{(θStar\ :\ EuclideanSpace\ ℝ\ (Fin\ dθ))\ (hMaximizer\ :\ IsMarginalMaximizer\ (U\ y)\ θStar)}
\leanline{4}{(initial\ :\ Lp\ (State\ N\ dθ\ dx)\ 2\ P)\ (noise\ :\ ℕ\ →\ Lp\ (State\ N\ dθ\ dx)\ 2\ P)}
\leanline{4}{(hGaussian\ :\ GaussianInnovations\ P\ initial\ noise)}
\leanline{4}{(k\ :\ ℕ)\ :}
\leanline{4}{Real.sqrt\ (∫\ z,\ ‖parameterCLM\ z\ -\ θStar‖\ \textasciicircum{}\ 2}
\leanline{6}{∂(laws\ P\ (U\ y)\ γ\ L\ hdiff\ hLip\ initial\ noise\ k))\ ≤}
\leanline{6}{Real.sqrt\ (2\ *\ (dθ\ :\ ℝ)\ /\ ((N\ :\ ℝ)\ *\ m))\ +}
\leanline{8}{(1\ -\ m\ *\ γ)\ \textasciicircum{}\ k\ *\ ULA.wasserstein2Real\ (P.map\ initial)\ (gibbsTarget\ (U\ y))\ +}
\leanline{8}{(33\ /\ 20\ :\ ℝ)\ *\ ((L\ :\ ℝ)\ /\ m)\ *}
\leanline{10}{Real.sqrt\ ((dx\ :\ ℝ)\ +\ (dθ\ :\ ℝ)\ /\ (N\ :\ ℝ))\ *\ Real.sqrt\ γ}
\end{leancode}
\leanlabel{How the Lean proof follows the argument}
\texttt{theorem1} receives the differentiable, strongly convex, smooth pair potential together with positive curvature, positive step, Gaussian independence and the marginal-maximizer premise. It splits on \texttt{hindex}, which records whether the full state has any coordinate. In the nonempty case \texttt{nonempty\_\allowbreak{}dimension\_\allowbreak{}bound} assembles scaled Gaussian transport and parameter concentration. In the empty case, \texttt{h\ensuremath{\theta}} and \texttt{hx} force both dimensions to zero, so the parameter error and every dimensional contribution vanish. The initial state is supplied as an \texttt{Lp} vector of exponent two; the Gaussian assumptions identify its algorithmic laws with the transition iterates.



The public theorem splits on whether the state coordinate index type is nonempty. In the nonempty case, instantiate the preceding theorem directly, retaining arbitrary independent \(d_\theta,d_x\). If the index type is empty, its parameter part forces \(d_\theta=0\); since \(N>0\), a positive \(d_x\) would supply a latent coordinate, so \(d_x=0\) as well. Normalize the Gibbs target by the same strong-convexity theorem and apply the singleton-state bound, where every error, transport distance and dimension term vanishes. This completes the displayed estimate in all dimensions. Existence and uniqueness of a permissible \(\theta_*\) are established separately by \texttt{exists\allowbreak{}Unique\_\allowbreak{}marginal\allowbreak{}Maximizer}, so the maximizer binder does not demand an unproved existence assumption.



Formal correspondence: theorem1 : Strongly convex IPLA for independent, arbitrary parameter and latent dimensions. All stationary-reference and concentration inputs are derived. The zero-dimensional cases are included without extra assumptions.
\end{leanbox}


Earlier results used here: \cite[Section~3.17.4]{LangevinLean2026ULASC}.

\begin{thebibliography}{ULA-LSI}
\bibitem[ULA-SC]{LangevinLean2026ULASC}
LangevinLean Project. \emph{Unadjusted Langevin under Strong Convexity}. LL-01, 2026. Draft report, content revision d50c0f64815b. \url{https://odakyildiz.github.io/LangevinLean/docs/pdf/UnadjustedLangevin/FullProof.pdf}.
\end{thebibliography}

\end{document}
